% TOP_PROBLEM: {"id": 3155, "title": "Domination in cubic graphs", "category_id": 3, "category": {"id": 3, "name": "graph_theory", "display_name": "Graph Theory", "description": "Problems involving graphs, networks, and their properties.", "slug": "graph-theory", "order_index": 3, "created_at": "2026-07-31T15:26:25.671Z"}, "status": "open", "problem_number": "OPG-37038", "set_id": 12}
% FIRST_POSTED: 2026-10-01
% ATTEMPT_STATUS: unresolved
% UnsolvedMath ID 3155; source catalogue code OPG-37038
% !TeX program = lualatex
% Research attempt by GPT-6 Astra Ultra; no claim of a new complete solution.
\documentclass[11pt,a4paper]{article}
\usepackage[margin=25mm]{geometry}
\usepackage{fontspec}
\setmainfont{lmroman10-regular.otf}[BoldFont=lmroman10-bold.otf,
 ItalicFont=lmroman10-italic.otf,BoldItalicFont=lmroman10-bolditalic.otf]
\usepackage{amsmath,amssymb,mathtools}
\usepackage{unicode-math}
\setmathfont{latinmodern-math.otf}
\AtBeginDocument{\let\setminus\smallsetminus}
\usepackage{xurl,hyperref}
\hypersetup{colorlinks=true,urlcolor=blue}
\setcounter{secnumdepth}{0}
\setlength{\parindent}{0pt}
\setlength{\parskip}{5pt}
\setlength{\emergencystretch}{3em}
\begin{document}
\section*{Domination in cubic graphs}\label{3155}
{\small UnsolvedMath ID 3155 · OPG-37038 · Graph Theory · OpenGarden}
\subsection{Definitions and mathematical statement}
Let $G=(V,E)$ be a finite simple undirected graph and write
$n=|V|$. For a vertex $v$, its closed neighbourhood is
$N[v]=\{v\}\cup\{u:uv\in E\}$. A set $D\subseteq V$ is dominating
if every vertex belongs to the closed neighbourhood of some vertex
of $D$, or equivalently
\[
 \bigcup_{v\in D}N[v]=V.
\]
The domination number $\gamma(G)$ is the minimum size of such a
set. Vertices in $D$ dominate themselves; there is no requirement
that they have another member of $D$ as a neighbour.

The graph is cubic if every vertex has degree three. It is
$3$-connected if $n\geq4$ and deleting any set of at most two
vertices leaves a connected graph. The proposed inequality is
\[
 \gamma(G)\leq\left\lceil\frac n3\right\rceil
 \qquad\text{for every $3$-connected cubic graph }G.
\]
Thus one asks for a set of at most the stated rounded-up size
meeting the closed neighbourhood of every vertex.

No independence or connectedness condition is imposed on $D$.
The graph need not be planar, bipartite, or have large girth.
Vertex connectivity is part of the hypothesis, so the problem
is not the assertion for arbitrary cubic graphs. The ceiling
has its ordinary integer meaning and matters at orders not
divisible by three. The quantification includes every admissible
finite order, rather than only an asymptotic estimate with an
unspecified additive error.
\subsection{Short English statement}
Can every 3-connected cubic graph be dominated by at most one third of its vertices, rounded up?
\subsection{Research attempt}
\paragraph{Scope and process.}
This is a GPT-6 Astra Ultra research attempt dated 1 October 2026, using
primary-source browsing and elementary mathematical checks. The canonical
definition above is retained. Results below are restricted results or
reductions, not a claim of a new solution to the entire catalogue question.
No formal proof assistant or external mathematical referee checked this text.


\paragraph{Literature and attack plan.}
Kelmans [R1] gives counterexamples of connectivity two to the older
assertion for all cubic graphs. Those counterexamples do not satisfy
the present three-connectivity hypothesis. This distinction is essential:
one cannot strengthen the induction hypothesis to arbitrary cubic
components after cutting a graph apart. We prove the desired inequality
for graphs with a suitable spanning cycle structure, retaining the exact
rounding cost, then examine why an arbitrary collection of cycles does
not suffice.

\paragraph{A direct construction from a Hamilton cycle.}
Suppose $G$ has a spanning cycle $v_0v_1\cdots v_{n-1}v_0$.
Take
\[
 D=\{v_0,v_3,v_6,\ldots,v_{3\lfloor(n-1)/3\rfloor}\}.
\]
This set has $\lceil n/3\rceil$ vertices. Consecutive chosen vertices
around the cycle are separated by at most three cycle edges, including
the final gap returning to $v_0$. In a gap of length three, its two
interior vertices are each adjacent to an endpoint; in shorter gaps
the same conclusion holds. Thus $D$ dominates the whole cycle, and
therefore all of $G$. Additional edges cannot destroy domination.
This proves the required upper bound for every Hamiltonian instance,
without needing regularity or three-connectivity in this restricted proof.

\paragraph{Keeping exact track of a spanning $2$-factor.}
More generally, suppose $F$ is a spanning union of vertex-disjoint
cycles of lengths $\ell_1,\ldots,\ell_s$. Apply the preceding choice
separately on each cycle and take the union. If $c_j$ counts cycles
with length congruent to $j$ modulo three for $j=1,2$, then
\[
 \gamma(G)\le\sum_{i=1}^s\left\lceil\frac{\ell_i}{3}\right\rceil
 =\frac n3+\frac{2c_1+c_2}{3}.
\]
This is a proved upper bound for each specific factor. In particular,
if $2c_1+c_2\le2$, its right side is $\lceil n/3\rceil$.
Indeed, it is an integer at least $n/3$ and strictly below $n/3+1$.
Thus a factor all of whose cycles have lengths divisible by three
works, as does one with only a single exceptional cycle, or with
two exceptional cycles both of length two modulo three.

The exact defect $2c_1+c_2$ matters. Bounding the number of cycles
without controlling their residues produces an additive loss, which
cannot be silently absorbed by the single ceiling in the question.
To prove the full target by this route one would need either a suitable
factor for every input graph, or a way to exploit edges between cycles
to save enough of these independently selected dominators.

\paragraph{Why a poor factor does not refute the target.}
In the cube graph $Q_3$, choose the two disjoint square faces obtained
by fixing one coordinate. They form a spanning $2$-factor with
$c_1=2$. The independent cycle construction gives four dominators,
whereas $\lceil8/3\rceil=3$. But in fact two opposite vertices,
$000$ and $111$, dominate the cube: every binary triple differs in
at most one coordinate from one of them. The edges joining the two
squares supply exactly the interactions discarded by treating the
cycles separately. Therefore a bad residue count for one factor is
an obstruction to this argument, not a counterexample to the conjecture.

There is also a universal elementary lower bound. In a cubic graph,
each closed neighbourhood has four vertices, so any dominating set
satisfies $4|D|\ge n$. Thus
$\gamma(G)\ge\lceil n/4\rceil$. The cube meets this bound. These
inequalities are compatible with the conjecture, but the counting lower
bound gives no selection mechanism at the upper threshold.

\paragraph{Verification.}
The script \texttt{scripts/check\_attempt\_batch02\_algebra\_2026\_10\_01.py}
constructs $K_4$, $K_{3,3}$, the triangular prism, the cube and the
Petersen graph. It checks that each is cubic and remains connected
after deletion of any set of at most two vertices. Exhaustive enumeration
of vertex subsets gives domination numbers $1,2,2,2,3$, respectively.
It also verifies the cyclic placement construction for cycle lengths
three through one hundred. These checks exercise both rounding classes
and the closed-neighbourhood convention. They do not enumerate all
three-connected cubic graphs of any unbounded order.

\paragraph{Final status and first remaining gap.}
\textbf{UNRESOLVED.} The Hamiltonian and stated $2$-factor subcases
satisfy the exact inequality. The missing step is a global choice or
exchange argument controlling the residue losses for every graph in
the hypothesis. Three-connectivity has not yet been converted into
such an argument, and lower-connectivity counterexamples cannot be
used to skip that requirement.

\subsection{Sources}
\begin{itemize}
\item[S1] ULAM.AI and the UnsolvedMath Contributors, supplied problems.json, record 3155 (OPG-37038), OpenGarden collection. Catalogue identity and source transcription; CC BY 4.0.
\url{https://huggingface.co/datasets/ulamai/UnsolvedMath}.
\item[S2] Open Problem Garden, Domination in cubic graphs. Original problem and discussion; the mathematical formulation below is expanded from the supplied source record.
\url{http://www.openproblemgarden.org/op/domination_in_cubic_graphs}.

\item[R1] Alexander Kelmans, \emph{Counterexamples to the Cubic Graph
Domination Conjecture}, 2006. Primary abstract checked 1 October 2026;
its counterexamples of connectivity two are distinguished from the
present three-connected question.
\url{https://arxiv.org/abs/math/0607512}.

\end{itemize}
\end{document}
